\documentclass[10pt,a4paper]{article}
\usepackage[utf8]{inputenc}\usepackage[T1]{fontenc}\usepackage[german]{babel}
\usepackage{graphicx,float,booktabs,amsmath,amssymb,siunitx,xcolor,tabularx,enumitem}
\usepackage[margin=2.0cm]{geometry}\usepackage[hidelinks]{hyperref}
\usepackage{newunicodechar}\newunicodechar{σ}{$\sigma$}\newunicodechar{–}{--}\newunicodechar{−}{$-$}\newunicodechar{≈}{$\approx$}
\usepackage{textcomp}
\sisetup{output-decimal-marker={,},group-separator={\,},group-minimum-digits=4}
\graphicspath{{bilder/}{../../docs/bilder/}}
\hypersetup{pdfauthor={Thomas Schmiegelt},pdftitle={Review-Paket Festigkeit: Spulentragwerk eines Stellarators}}
\newcommand{\bild}[3]{\begin{figure}[H]\centering\includegraphics[width=#2\linewidth]{#1}\caption{#3}\end{figure}}
\newcommand{\lesehilfe}[1]{\begin{quote}\small\textit{So liest man das Bild:} #1\end{quote}}
\newcommand{\kasten}[2]{\par\smallskip\noindent\fcolorbox{gray}{gray!8}{\parbox{0.96\linewidth}{\textbf{#1} #2}}\par\smallskip}
\title{\textbf{Review-Paket Festigkeit}\\[2pt]\large Spulentragwerk eines Stellarators -- Modell, Ergebnisse und
acht Fragen an Festigkeitsfachleute}
\author{Thomas Schmiegelt \textbar{} Hobbyprojekt \texttt{stellarator-design}}
\date{3.~November 2026}
\makeatletter\def\ps@kifuss{\def\@oddhead{}\def\@evenhead{}\def\@oddfoot{\hfil{\footnotesize\thepage\quad\textperiodcentered\quad KI-generiert: Text, Programmcode und Abbildungen von KI-Agenten; Ziel und Richtung vom Menschen}\hfil}\let\@evenfoot\@oddfoot}\let\ps@plain\ps@kifuss\makeatother
\begin{document}\maketitle\pagestyle{kifuss}
\begin{center}\fbox{\parbox{0.96\textwidth}{\small \textbf{KI-generierter Bericht.} Text, Programmcode und Abbildungen dieses Berichts wurden von KI-Agenten erzeugt. Gerechnet haben die genannten Löser. Der Mensch (Thomas Schmiegelt) gab Ziel und Richtung vor, entschied an den Gabelungen und prüfte Bilder und Ergebnisse auf Plausibilität.}}\end{center}


\kasten{Hobbyprojekt, nicht kommerziell.}{Keine Verbindung zu einem Kundenprojekt, Forschungsinstitut oder
Unternehmen, nicht begutachtet. Rechnungen, Programmcode und Textentwurf entstanden mit Unterstützung von
KI-Agenten; der Mensch gab Ziel und Richtung vor. Alle Zahlen stammen aus vereinfachten Modellen (linear-elastisch,
starre Lager, nur Magnetkraft). Gerade deshalb bitte ich um euren kritischen Blick: Die Fragen stehen in
Abschnitt~\ref{sec:fragen}.}

\noindent\textbf{Für wen.} Festigkeitsfachleute (FEM, Regelwerk) ohne Fusionshintergrund. Alles Fusionsspezifische
ist kurz erklärt. Lesezeit etwa 20 Minuten; für die Fragen allein genügen Abschnitt~\ref{sec:kurz} und
\ref{sec:fragen}.

\section{Auf einen Blick}\label{sec:kurz}
\begin{itemize}[itemsep=1pt]
\item \textbf{Bauteil:} Gehäuse aus Stahl 1.4429 (316LN) um 24 supraleitende Spulen, verbunden durch Stahlverbinder,
gelagert auf einem inneren und äußeren Stützring. Last: Lorentzkraft im Magnetbetrieb, bis \SI{1,1}{GN} je Spule.
\item \textbf{Bewertung:} σ99 = 99-\%-Wert der von-Mises-Spannung über das Stahlvolumen; Grenze \SI{500}{MPa}
($\tfrac23 R_{p0,2}$ bei \SI{4}{K}, ITER-Güte C3).
\item \textbf{Ergebnis:} Von 199 ausgewerteten Entwürfen einer Studie halten 7 alle Grenzen ein. Empfohlen: Entwurf 0169 mit
σ99 = \SI{308}{MPa}, auch mit feinerem Netz (\SI{319}{MPa}).
\item \textbf{Hauptproblem:} Die Höchstwerte sitzen an den Anschlüssen der Verbinder und wachsen mit jeder
Netzverfeinerung. Bei vier von fünf Kandidaten steigt σ99 mit feinerem Netz um 34--58~\%.
\item \textbf{Neu (3.11.):} Bei 0169 lagen die Verbinder so dicht, dass sie fast eine Platte bildeten. Nachgerechnet:
Getrennte Stege sind netzstabil (+1~\%), eine lückenlose Platte nicht (+47~\%).
\end{itemize}

\section{Worum es geht}
Ein Stellarator hält ein über 100 Millionen Grad heißes Plasma in einem verdrillten Magnetfeld. Das Feld kommt aus
supraleitenden Spulen mit unregelmäßiger, dreidimensionaler Form. Der Entwurf hier: \SI{1}{GW} elektrische
Nettoleistung, großer Radius $R_0=\SI{10,4}{m}$, Feld $B_0=\SI{6,5}{T}$, Hochtemperatur-Supraleiter (REBCO), Betrieb
bei etwa \SIrange{4}{20}{K}. Jede Spule trägt rund \SI{20}{MA}. Im eigenen und im Feld der Nachbarn wirkt auf sie die
Lorentzkraft $f=J\times B$: \SIrange{50}{140}{MN/m} entlang der Spule, in Summe bis \SI{1,1}{GN} je Spule zur
Anlagenmitte. Gesucht ist ein Tragwerk aus Spulengehäusen, Verbindern und Lagern, das diese Kräfte aufnimmt, ohne
dass der Supraleiter mehr als \SI{0,3}{\%} gedehnt wird.

\bild{modell3d_0169_schnitt}{0.72}{Entwurf 0169 im Treppenschnitt (3D-Modell in PicoGK).}
\lesehilfe{Von innen nach außen: Plasma (rosa), Erstwand und Blanket (gelb), Abschirmung (grün), Vakuumgefäß (grau),
dann die 24 Spulengehäuse aus Stahl (blau) mit dem Wickelpaket aus Supraleiter (orange). Zwischen den Gehäusen sitzen
die Verbinder (dunkelblau). Stützzylinder innen und Ring außen sind nur gezeichnet: Sie zeigen, wohin die Lagerkräfte
gehen, sind aber weder berechnet noch ausgelegt.}

\section{Modell}
\begin{description}[itemsep=1pt,style=nextline]
\item[Geometrie] Spulenmittellinien aus der Optimierung (simsopt \cite{simsopt}). Je Spule ein Hexaedernetz entlang
der Bahn: 64 Stationen bei 24 Spulen, im Querschnitt 2 Wandzellen, \SI{2}{cm} Verguss und $4\times4$ Zellen
Wickelpaket. Verbinder als Hexaederblöcke zwischen den Gehäusen.
\item[Werkstoffe] Gehäuse 1.4429 (X2CrNiMoN17-13-3, AISI 316LN), $E=\SI{200}{GPa}$. Wickelpaket homogenisiert.
Verguss Quarzsand-Epoxid, $E=\SI{10}{GPa}$. Verbinder Stahl, $E$ als Stellgröße \SIrange{20}{200}{GPa}.
\item[Last] Lorentzkraftdichte aus allen Spulen (Biot-Savart), nur Magnetbetrieb.
\item[Lager] Flecken auf der Gehäuseaußenseite am kleinsten und größten Radius, \textbf{starr eingespannt}.
\item[Löser] Elmer \cite{elmer}, linear-elastisch, lineare 8-Knoten-Hexaeder (Elmer-Typ 808), MUMPS.
\item[Kette] Plasma $\to$ Spulen $\to$ Package $\to$ Festigkeit $\to$ Temperatur, 1--3 Minuten je Entwurf. Eine
Studie mit 200 Entwürfen variierte Spulenzahl, Spulenabstand, Verbinderzahl und -steifigkeit, Wanddicken und
Lagerradius.
\end{description}

\section{Wie die Varianten entstehen: automatisch gerechnet, automatisch verworfen}\label{sec:auto}
Kein Entwurf wurde von Hand gezeichnet. Eine Kette von Programmen erzeugt aus elf Stellgrößen Spulen, Package,
Festigkeitsmodell, Temperaturen und ein 3D-Modell und bewertet das Ergebnis gegen acht Grenzen. Die Stellgrößen
verteilt ein Latin-Hypercube-Plan gleichmäßig über den Bereich. Die KI schrieb die Kette und startete sie; die 200
Entwürfe liefen ohne Eingriff in 8,5 Stunden durch.

\bild{ablauf}{0.98}{Die automatische Kette der Entwurfsstudie und was davon übrig bleibt.}
\lesehilfe{Oben die Stationen, die jeder Entwurf durchläuft. Links unten der Trichter: Von 200 geplanten Entwürfen
scheiterte bei einem der Löser (Abbruch ohne Ergebnis), bei 11 durchdringen sich die Gehäuse (nicht baubar). Von den übrigen 188
halten 7 alle Grenzen ein, 3 weitere verfehlen eine knapp. Rechts unten, woran die 181 verworfenen scheitern: an der
Grenze mit der kleinsten Reserve. Meist ist es die Spannung im Stahl, aber fast die Hälfte scheitert an anderem
(Feldgenauigkeit, Dehnung des Supraleiters, Neutronenbelastung). Ein reiner Festigkeitsentwurf würde sie nicht
aussortieren.}

\bild{galerie}{0.98}{16 verworfene Entwürfe, gleichmäßig über die Rangfolge verteilt (3D-Modelle in PicoGK).}
\lesehilfe{Oben links der Entwurf, der am weitesten von den Grenzen entfernt ist, unten rechts der, der am knappsten
scheitert. Über jedem Bild: Nummer, Spulenzahl, σ99 und die Grenze, an der er scheitert. „Reserve“ ist der Abstand
zur Grenze relativ zur Grenze ($-100$~\% heißt doppelt so viel wie zulässig). Roter Rahmen: deutlich verfehlt, gelb:
knapp. Hohe Spannung allein entscheidet nicht: 0021 hat nur \SI{558}{MPa}, aber sein Supraleiter bekäme ein
Vielfaches der zulässigen Neutronenbelastung. Im Katalog des Variantenberichts sind alle 188 baubaren Entwürfe so
abgebildet.}

\bild{film_feld}{0.82}{Bild aus dem Film aller Varianten (Sekunde 30): Die verworfenen versinken.}
\lesehilfe{Jeder Ring ist ein ganzer Entwurf aus der Studie. Rot: deutlich verfehlt, gelb: knapp, blau: zulässig.
Die unzulässigen versinken der Reihe nach, die schlechtesten zuerst; am Ende bleiben die sieben blauen. Der Film
liegt im Repository (\texttt{docs/filme/varianten.mp4}).}

\section{Wie das Tragwerk entstanden ist}
Das Tragwerk wurde nicht vorgegeben, sondern aus einer Kraftanalyse abgeleitet. Das erste Modell kam auf
\SI{2129}{MPa}. Statt Parameter zu verändern, wurde zuerst untersucht, wie die Kraft entsteht.
\begin{enumerate}[itemsep=1pt]
\item \textbf{Kraft entlang der Spule:} \SIrange{50}{140}{MN/m}, Summe je Spule bis \SI{1,1}{GN} nach innen.
\item \textbf{Seil oder Balken:} Ein gekrümmtes Stück trägt Querlast als Zug ($T=f_N/\kappa$), ein gerades wird
gebogen. Die größte Last (\SI{\approx130}{MN/m}) sitzt auf dem fast geraden inneren Schenkel ($\kappa<\SI{0,08}{1/m}$
über rund \SI{7}{m}).
\item \textbf{Handrechnung:} Balken auf zwei Stützen, $q\approx\SI{130}{MN/m}$, $l\approx\SI{7}{m}$:
$M=ql^2/8\approx\SI{800}{MNm}$; Hohlquerschnitt $H=\SI{1,36}{m}$, $h=\SI{0,755}{m}$, $W\approx\SI{0,38}{m^3}$,
also $\sigma_B\approx\SI{2100}{MPa}$. Die FEM ergab \SI{2129}{MPa}: Die Spannung kam aus der Biegung des inneren
Schenkels.
\item \textbf{Fehler im Ausgangsmodell:} Gelagert war nur außen, die Kraft zieht nach innen.
\item \textbf{Vorbilder:} Beim Tokamak liegt der gerade Schenkel an einem Stützzylinder an (Druck statt Biegung); bei
ARIES-CS gibt eine durchgehende Schale die Querlast von Spule zu Spule weiter.
\end{enumerate}

\bild{festigkeit_tragwerk}{0.95}{Gleiche 16 Spulen, fünf Tragwerke A--E.}
\lesehilfe{Oben vier Kennwerte je Tragwerk, unten die Spannung von oben gesehen (blau niedrig, rot hoch). A: Lager
außen, 6 Verbinder je Paar. B: Lager innen ($-32$~\%, erwartet waren $-30$ bis $-50$~\%). C: innen und außen. D: 12
steife Verbinder je Paar (kürzere freie Biegelänge, $M\propto l^2$). E: zusätzlich dickste Wände. Die Spannung sinkt
von \SI{2129}{} auf \SI{481}{MPa}, der Stahl wächst von \SI{15,6}{} auf \SI{37}{kt}. Seitdem ist „innen und außen“
Standard; Verbinder, Wände und Lagerradius sind Stellgrößen der Studie.}

\section{Bewertung und Grenzen}\label{sec:grenzen}
\begin{table}[H]\centering\small\begin{tabular}{@{}lll@{}}\toprule
Größe & Grenze & Herkunft \\\midrule
σ99 im Gehäusestahl & \SI{500}{MPa} & $S_m=\tfrac23R_{p0,2}$ bei \SI{4}{K}, ITER-Güte C3 ($R_{p0,2}\ge\SI{750}{MPa}$) \cite{mitchell}\\
Membranspannung $P_m$ je Wandschnitt & \SI{500}{MPa} & Linearisierung über die Wanddicke \\
$P_m+P_b$ & \SI{750}{MPa} & $1{,}5\,S_m$ \\
Dehnung Wickelpaket & \SI{0,3}{\%} & REBCO irreversibel ab \SIrange{0,45}{0,72}{\%} \\\bottomrule
\end{tabular}\end{table}

\bild{stahl_kennwerte}{0.78}{Festigkeit von 1.4429 / 316LN bei Raumtemperatur und bei \SI{4}{K}.}
\lesehilfe{Blau die Streckgrenze, grün die Zugfestigkeit, gelb die daraus folgende zulässige Spannung. Bei
Raumtemperatur (links, Tabellenwert \cite{schweizerfn}) wären nur rund \SI{200}{MPa} zulässig. Bei \SI{4,2}{K}
gemessen ist die Streckgrenze mehr als dreimal so hoch (\SI{1070}{MPa} \cite{fedorikova2023}; mit hohem Stickstoff
in dicken Blechen über \SI{1000}{MPa} \cite{sakurai2024}). Die gestrichelte Linie ist die verwendete Grenze. Das
Gehäuse ist im Betrieb kalt; beim Abkühlen fließt noch kein Strom.}

\section{Ergebnisse}
\bild{doe1_uebersicht}{0.85}{Entwurfsstudie: 199 Entwürfe, σ99 gegen sechs Stellgrößen.}
\lesehilfe{Jeder Kreis ist ein Entwurf, die Farbe gibt die Spulenzahl an. Gefüllte Kreise halten alle Grenzen ein.
Die gestrichelte Linie ist die \SI{500}{MPa}-Grenze, die y-Achse ist logarithmisch. Nur Entwürfe mit 20 oder 24
Spulen kommen darunter. Den deutlichsten Trend zeigt der Abstand Spule–Plasma (unten Mitte): Liegen die Spulen
weiter weg, sind sie stärker verformt.}

\paragraph{Netzstudie.} Die Kandidaten wurden mit 1,5-mal so vielen Stationen entlang der Spule und mit 3 statt 2
Wandzellen nachgerechnet.
\bild{doe1_netz}{0.95}{Netzstudie: gleiche Spulen und Parameter, anderes Netz.}
\lesehilfe{(b) zeigt σ99, (c) den Höchstwert. Eine feinere Auflösung der Wanddicke („Querschnitt fein“) ändert kaum
etwas. Mehr Stationen entlang der Spule („1,5× Stationen“) erhöhen σ99 bei vier von fünf Kandidaten um 34--58~\%,
bei 0169 nur um 4~\%. Der Höchstwert steigt bei allen, er hängt am Netz.}

\paragraph{Ort der Spitzen.} 77--100~\% der 200 höchstbelasteten Elemente liegen weniger als \SI{1}{m} von einem
Verbinderanschluss entfernt. Bei 0169 steigt der Höchstwert mit dem feineren Netz von \SI{717}{} auf
\SI{1262}{MPa}, σ99 nur von \SI{308}{} auf \SI{319}{MPa}.

\paragraph{Kategorien.} Die Spannung wurde über jeden Wandschnitt linearisiert, in Membran ($P_m$) und Membran
plus Biegung ($P_m+P_b$).
\bild{doe1_kategorien}{0.95}{Spannungskategorien der zulässigen Entwürfe.}
\lesehilfe{(a) Höchstwerte: Sie liegen an der Einzelspitze am Verbinderende (örtliche Membranspannung $P_L$). (b)
Wert, den 99~\% der Wandschnitte einhalten: Bei den besten Entwürfen bleibt $P_m$ unter \SI{400}{MPa}. E (test) ist
der alte Tragwerksversuch E zum Vergleich.}

\paragraph{Verbinder als Stege oder als Platte.} Bei 0169 (13 Verbinder je Paar) lagen benachbarte Verbinder genau
eine Verbinderbreite auseinander. Sie berührten oder überlappten sich, und in den Überlappungen wurde Stahl doppelt
gezählt. Vermutung vor der Rechnung: Die fast durchgehende Platte macht 0169 netzstabil. Nachgerechnet wurden drei
Varianten mit gleichen Spulen.
\bild{verbinder_modi}{0.72}{Entwurf 0169: drei Arten der Verbinder, Standard- und feines Netz.}
\lesehilfe{Blau das Standardnetz, orange 1,5-mal so viele Stationen. Die Zahl über jedem Paar ist der Zuwachs durch
das feinere Netz. Getrennte Stege (Abstand mindestens eine Breite) liegen höher (\SI{387}{MPa}), bleiben aber stabil
(+1~\%). Eine lückenlose Platte mit gemeinsamen Knoten springt um 47~\% auf fast \SI{500}{MPa}. Die Vermutung ist
widerlegt: Die Netzstabilität von 0169 kommt nicht von der Platte. Die Spitzen liegen in gut geformten Elementen
(skalierte Jacobi-Qualität 0,92--0,96), zu 86--96~\% im Gehäuse. Auch mit getrennten Stegen bleibt 0169
zulässig.}

\paragraph{Weichere Verbinder.} Werden alle Verbinder von 0169 gleichmäßig weicher, steigt σ99 (+16~\% bei
\SI{50}{GPa}). Die Verbinder sind der Entlastungsweg.

\section{Bekannte Schwächen}
\begin{enumerate}[itemsep=1pt]
\item Lager starr eingespannt; der Stützzylinder fehlt als Bauteil (optimistisch).
\item Verbinderanschlüsse scharfkantig, daher netzabhängige Spitzen. Die Bewertung gilt für die Wände, nicht für die
Anschlüsse.
\item Nur Primärlast: keine Abkühlung von \SI{293}{} auf \SI{4}{K}, kein Eigengewicht, keine Ermüdung, keine
Bruchmechanik.
\item Wickelpaket homogenisiert, Verguss ideal verbunden (keine Ablösung).
\item σ99 ist keine Regelwerksgröße, sondern eine Kenngröße für die Entwurfssuche.
\item Lineare Hexaeder mit nur 2 Zellen über die Wanddicke.
\end{enumerate}

\section{Fragen an euch}\label{sec:fragen}
\begin{enumerate}[itemsep=3pt]
\item \textbf{Bewertungsgröße:} Ist σ99 als Suchkriterium vertretbar, oder sollten wir von Anfang an mit
linearisierten Kategorien ($P_m$, $P_L$, $P_b$) entlang ausgewählter Schnitte bewerten?
\item \textbf{Spitzen am Anschluss:} Die Höchstwerte am Verbinderende wachsen mit dem Netz (0169: \SI{717}{} $\to$
\SI{1262}{MPa}). Wie behandelt ihr das üblicherweise: Strukturspannung nach dem Hot-Spot-Verfahren, Kerbspannung mit
Ausrundungsradius, Submodell?
\item \textbf{Anschlussform:} Welche Form würdet ihr für einen Stahlverbinder zwischen zwei Gehäusen wählen (Spalt
etwa \SIrange{0,2}{1}{m}, Kräfte im MN-Bereich je Verbinder): Ausrundung, auslaufende Dicke, längere
Anschlussfläche, durchgehende Schale?
\item \textbf{Platte oder Stege:} Im Modell ist die lückenlose Platte netzempfindlich (+47~\%), getrennte Stege sind
es nicht (+1~\%). Ist das ein bekannter Effekt (Steifigkeitssprung am Plattenrand)? Spricht etwas gegen eine
durchgehende Schale (Montage, Wärmedehnung beim Abkühlen, Prüfbarkeit der Schweißnähte)?
\item \textbf{Lastfälle:} Welche fehlen, in welcher Reihenfolge? Vorschlag: Abkühlung, Eigengewicht,
Schnellabschaltung (Quench), zyklische Last für Ermüdung.
\item \textbf{Lagerung:} Wie würdet ihr den Stützzylinder abbilden, als eigenes Bauteil mit Kontakt oder als
Federrandbedingung? Welche Steifigkeit wäre realistisch?
\item \textbf{Werkstoff bei 4~K:} Reicht $\tfrac23R_{p0,2}$ (ITER-Magnetkriterien) als Grenze, oder ist bei \SI{4}{K}
die Bruchmechanik maßgebend ($K_{IC}$, Rissfortschritt)?
\item \textbf{Netz:} Wie stark unterschätzen lineare 8-Knoten-Hexaeder mit 2 Zellen über die Wanddicke die
Biegespannung (Schubversteifung)? Lohnen quadratische Elemente (20 Knoten), oder reicht eine dritte Zelle? Unser
Versuch mit 3 Wandzellen änderte σ99 kaum.
\end{enumerate}

\section*{Unterlagen}
Projektbericht (Kapitel Festigkeit), Variantenbericht mit Steckbriefen aller zulässigen Entwürfe und Katalog aller
188 baubaren Entwürfe, Quellcode (\texttt{stellarator/kopplung.py}, \texttt{hexmodell.py}, \texttt{doe\_auswertung.py})
und ein Film aller Varianten: \url{https://github.com/ThomasSchmiegelt/stellarator-design} (MIT-Lizenz).

\begin{thebibliography}{9}\small
\bibitem{simsopt} M.~Landreman u.\,a.: SIMSOPT: A flexible framework for stellarator optimization. J.\ Open Source
  Softw.\ 6 (2021) 3525. doi:10.21105/joss.03525.
\bibitem{elmer} Elmer FEM, CSC -- IT Center for Science, \url{https://www.elmerfem.org}.
\bibitem{mitchell} N.~Mitchell u.\,a.: ITER-TF-Struktur, Güten C1--C3 (IAEA-FEC-Beitrag IT/1-4); Jahrgang und
  Wortlaut nicht geprüft.
\bibitem{schweizerfn} Festigkeitswerte Stahl, nach Roloff/Matek, Maschinenelemente; schweizer-fn.de,
  \url{https://www.schweizer-fn.de/festigkeit/festigkeitswerte/stahl/stahl_start.php} (abgerufen 3.11.2026).
\bibitem{fedorikova2023} A.~Fedoriková, P.~Petroušek, T.~Kvačkaj, R.~Kočiško, M.~Zemko: Development of Mechanical
  Properties of Stainless Steel 316LN-IG after Cryo-Plastic Deformation. Materials 16 (2023) 6473.
  doi:10.3390/ma16196473.
\bibitem{sakurai2024} T.~Sakurai, O.~Umezawa, Y.~Ono: Strength and Fracture Toughness of Type 304 and 316 Austenitic
  Stainless Steels at 4.2~K. IOP Conf.\ Ser.: Mater.\ Sci.\ Eng.\ 1302 (2024) 012002.
  doi:10.1088/1757-899X/1302/1/012002.
\end{thebibliography}
\end{document}
